\section{Setting}\label{sec:setting}

\subsection{Packages, quotients and residuals}\label{sec:setting-apparatus}

\begin{definition}[Theory package]\label{def:theory-package}
A \emph{theory package} $\mathcal T = (Z, f, \Sigma_f, E, \mathcal A)$ consists of a carrier $Z$ of
states; a \emph{lens} $f$, the map that says what an observer of the package reads off a state; its
\emph{definability structure} $\Sigma_f$, the family of distinctions expressible through $f$; a
\emph{completion} $E$, a map on descriptions that packages a description into a finished object; and
an \emph{audit} $\mathcal A$, the checks the package applies to its own records. It comes with a
transition kernel $K$ on $Z$ whose support $\supp K \subseteq Z \times Z$ lists the allowed one-step
moves. The package is \emph{formed} when its completion has a declared fixed point with its records
and audit included in the package \citep{Tsiokos2026FoundI}.
\end{definition}

The carrier is finite in every application of this paper; in Lean the carrier is an arbitrary type
and finiteness is imposed only where a statement uses it. Formedness is a proposition supplied with
the package. None of the results below uses the internal structure of $f$, $\Sigma_f$ or $E$.

\paragraph{Quotients and split pairs.} A \emph{quotient} is a map $q : X \to Q$; it records which
states an observer can distinguish. Say $q$ \emph{refines} $r$, written $q \preceq r$, when
$q(x) = q(x')$ implies $r(x) = r(x')$. For an update $F : X \to Y$ and a readout $r : Y \to R$, the
\emph{split-pair obstruction} is
\begin{equation}\label{eq:split}
  \Delta(q, F, r) = \{(x, x') : q(x) = q(x'),\ r(F(x)) \ne r(F(x'))\}.
\end{equation}
The readout descends through $q$ exactly when $\Delta(q,F,r) = \varnothing$, and a nonempty
obstruction is repaired by refining the source or coarsening the target \citep{Tsiokos2026FoundIV}.

\paragraph{Closure deficit and adequacy residual.} On a history carrier, let $X_t$ be the history up to
time $t$, let $Y^+_{t,\tau}$ be the declared future observables up to horizon $\tau$, and for a quotient
$\Pi$ of histories let $Y^\Pi_t$ be the class of $X_t$ under $\Pi$. The \emph{closure deficit}
$\CD_\tau(\Pi) = I(X_t; Y^+_{t,\tau} \mid Y^\Pi_t)$ is the information about the future that the
history holds beyond what $\Pi$ keeps \citep{Tsiokos2026FoundIII}. It is used for a current quotient
$Q$ and a predictive quotient $M$ refining it. For a covariance
operator $C$, a native probe family $L$ and a dissolving family $D$, the \emph{adequacy residual} is
the Schur complement
\begin{equation}\label{eq:xi}
  \Xi_C(D \mid L) = K_{DD} - K_{DL} K_{LL}^{\dagger} K_{LD},
\end{equation}
and adjoining a probe $M$ changes it by the chain rule
\begin{equation}\label{eq:xi-chain}
  \Xi_C(D \mid [L; M]) = \Xi_C(D \mid L) - K_{DM\mid L}\, K_{MM\mid L}^{\dagger}\, K_{MD\mid L}
\end{equation}
\citep{Tsiokos2026Adequacy}. Both identities are imported.

\paragraph{Structural roles.} Repairs are typed by the structural roles of \FoundI{}:
P1 lower-operator rewrite, P2 support gating, P3 route reconciliation, P4 staging and refinement,
P5 packaging and completion, and P6 audit and accounting. A \emph{repair sort} is one of P1--P5.
P6 governs the other moves and is never itself emitted as a repair.

\subsection{Laws, general theorems and certificate specifications}\label{sec:setting-status}

Every entry E1--E16 has a general Lean statement and receives one of four statuses.
\begin{itemize}
  \item A \emph{law} is a general statement together with an instance: a concrete system for which
    Lean proves the hard hypothesis, so that the conclusion is established for that system. An
    instance is \emph{weak} when some of its data are declared rather than derived --- for example
    a completeness proposition declared to be true, or a record sort with a single element.
  \item A \emph{general theorem} has ordinary hypotheses that can be checked system by system, and
    nothing is left open.
  \item A \emph{certificate specification}, printed as a Template, is an implication $H \Rightarrow C$
    in which $H$ is a named certificate that carries the difficult content, and supplying $H$ for a
    real system is the \emph{open obligation}. A \emph{conditional classification} is a certificate
    specification whose conclusion is that exactly one of several statuses holds: given a complete
    status record, which already asserts that one branch is valid and that every record for the same
    claim has the same status, the statuses are mutually exclusive. The content is the exclusivity of
    the branch conditions and the fields each branch requires; the existence of a valid branch is
    part of the certificate.
  \item A \emph{definition with consistency lemmas} is an entry whose principal statements follow by
    unfolding a definition or projecting its fields.
\end{itemize}

Three words recur in the classifications. A \emph{status record} is a carried record that assigns
one status, or \emph{branch}, of a classification to a claim; it is \emph{branch-valid} when the
conditions listed for that branch hold for the claim. A residual or defect is \emph{statused} when
the ledger carries a record that names it and assigns it a status (unresolved, discounted, charged
to a budget line), rather than leaving it unaccounted for.

\subsection{The definitional layer}\label{sec:setting-defs}

\begin{definition}[D1: repair join]\label{def:d1}
For quotients $q : X \to A$ and $r : X \to B$ on a common carrier, the \emph{repair join} is
$(q \vee r)(x) = (q(x), r(x))$. A \emph{source repair} by $r$ replaces $q$ by $q \vee r$.
\end{definition}

\begin{proposition}[repair-join semilattice]\label{prop:d1}
For quotients $q, r, s, t$ on $X$:
\begin{enumerate}
  \item $q \vee r \preceq q$ and $q \vee r \preceq r$;
  \item if $t \preceq q$ and $t \preceq r$, then $t \preceq q \vee r$;
  \item $q \vee q$, $r \vee q$ and $q \vee (r \vee s)$ have the same fibres as $q$, $q \vee r$ and
    $(q \vee r) \vee s$ respectively;
  \item replacing $q$ or $r$ by a quotient with the same fibres does not change the fibres of
    $q \vee r$.
\end{enumerate}
Thus, up to equality of fibres, $\vee$ is the greatest lower bound in the refinement order, and it
is idempotent, commutative and associative.
\end{proposition}

\begin{proof}
All four parts come from the fact that the fibre relation of the join is the conjunction of the two fibre relations.
Projection onto the two coordinates gives (1). If $t$ refines $q$ and $r$, then $t(x) = t(x')$
gives both $q(x) = q(x')$ and $r(x) = r(x')$, hence $(q \vee r)(x) = (q \vee r)(x')$. The fibre
relation of $q \vee r$ is the conjunction of the fibre relations of $q$ and $r$, which gives (3)
and (4).
\end{proof}

The join supplies a common refinement. It does not show that a repair occurs, that the system
produced it, or that it reduces an obstruction; those are the obligations of E1.

\begin{definition}[D2: predictive surplus]\label{def:d2}
For a current quotient $Q$, a predictive baseline $M$ and a horizon $\tau$ with
$\CD_\tau(M) \le \CD_\tau(Q)$, the \emph{predictive surplus} is
$S_\tau(Q \Vert M) = \CD_\tau(Q) - \CD_\tau(M)$. In the operator setting the surplus is measured by
the trace $\tr \Xi_C(D \mid L)$.
\end{definition}

\begin{proposition}[conditional surplus inequalities]\label{prop:d2}
\begin{enumerate}
  \item For rational scores $d_Q, d_M$, the difference $d_Q - d_M$ is nonnegative if and only if
    $d_M \le d_Q$; it is zero when $d_M = d_Q$, and equals $d_Q$ when $d_M = 0$.
  \item If $\Xi_C(D \mid L)$ is positive semidefinite, its trace is nonnegative.
  \item Suppose the residual $\Xi'$ after adjoining $M$ satisfies the chain rule
    \eqref{eq:xi-chain} and the subtracted term is positive semidefinite. Then
    $\tr \Xi' \le \tr \Xi_C(D \mid L)$.
  \item If moreover $M = BL$ for some matrix $B$ and $K_{DM\mid L} = 0$, then
    $\tr \Xi' = \tr \Xi_C(D \mid L)$.
\end{enumerate}
\end{proposition}

\begin{proof}
(1) is arithmetic. (2) and (3) use that the trace is nonnegative and monotone on the positive
semidefinite cone; in (3) the chain rule writes $\Xi'$ as $\Xi_C(D\mid L)$ minus a positive
semidefinite term. In (4) the subtracted term vanishes because its left factor is zero.
\end{proof}

The proposition is conditional: the score comparison, the chain-rule identity and positive
semidefiniteness are hypotheses. Refining a quotient does not by itself make the surplus
nonnegative, and scores taken at different horizons or against different future families are not
compared.

The central notion of the paper is a record that the system itself holds. Informally, a record is
the system's own when it is a coordinate of the system's state, reached along a legitimate trajectory
of the system's own kernel within scope, tagged as coming from committed state or from audited
records rather than from an outside source, and flagged as generated by the system within scope.

\begin{definition}[D3: carried record]\label{def:d3}
Let $\mathcal T$ have carrier $Z$ and kernel support $\supp K$. A record value $\rho$ with source tag
$\sigma$, generation flag $g$ and scope flag $i$ is \emph{carried} if
\begin{enumerate}
  \item $\rho = \rho_{\mathrm{of}}(z)$ for a declared coordinate map $\rho_{\mathrm{of}} : Z \to R$
    and some state $z$ of a declared trajectory $\tau$;
  \item $\tau$ starts at a legitimate start and every step up to $z$ is in scope and lies in
    $\supp K$; and
  \item $\sigma$ is \emph{committed state} or \emph{audited cell records}, and $g = i = $ true.
\end{enumerate}
A \emph{carried instrument} is an instrument whose visibility, threshold and check-rule records are
all carried and whose record list is declared complete.
\end{definition}

A value from a fallback, a simulation trace, an ablation record or an independent witness is thus
not the system's own record, whatever flags it carries (\Cref{lem:d3}).

\begin{definition}[D4: E-system and lawful repair step]\label{def:d4}
An \emph{E-system} $\mathbb S = (\mathcal T, I_S, \Lambda_S, \mathfrak R_S)$ consists of a formed
package $\mathcal T$; a carried instrument $I_S$ with predicates \emph{detects}, \emph{gate allows}
and \emph{re-audits}; a carried ledger $\Lambda_S$, a list of entries declared complete, each
carried; a repair generator $\mathfrak R_S$ sending each defect record to a repair move of some sort
in P1--P5 with a carried move record and a budget line; and an admissibility predicate for moves.
\end{definition}

\begin{definition}[lawful repair step]\label{def:d4-step}
For states $z, z'$, a defect record $d$ and an audit record $a$, a \emph{lawful repair step}
$z \to z'$ requires: $d$ is carried; $I_S$ detects $d$ at $z$; the gate allows $\mathfrak R_S(d)$ at
$z$; the move is admissible at $z$; its budget line is an entry of $\Lambda_S$; $(z, z') \in \supp K$;
$I_S$ re-audits $z'$ with record $a$; and $a$ is carried.
\end{definition}

Each of detection, gate permission, kernel transition and re-audit is necessary (\Cref{lem:d4}).

The step is a lawful possible transition with linked records. It does not show that the transition
takes place along a trajectory, nor recurrence, reduction of an obstruction, feasibility of a family
of repairs within the budget, or maintenance over time.

\begin{definition}[D5: closed-loop scope]\label{def:d5}
A \emph{challenge history} lists challenge episodes and repair-typed audit entries (defects, moves,
audits), each entry associated with an episode, together with a proposition declaring the history
complete. It is \emph{closed-loop} for $\mathbb S$ when it is declared complete and every
repair-typed entry in it is carried.
\end{definition}


A single listed entry with source tag fallback, or a listed entry with a false generation or scope
flag, defeats closed-loop scope (\Cref{lem:d5}). A history with no
repair entries is closed-loop. This is intended: closed-loop scope rules out
external repair; it does not assert that any repair took place.

\begin{definition}[D6: probe economy]\label{def:d6}
A \emph{probe economy} has a finite catalog of probes, a carried active family $L$ with support and
weights, an exposure budget carried in $\Lambda_S$, and three moves: \emph{allocation}, which
reweights the current support; \emph{acquisition} of a catalog probe $M$ outside the support; and
\emph{retirement} of an active probe. An acquisition is \emph{strict} when $M$ is not same-family
saturated with respect to $L$; in the operator setting, $M = BL$ with $K_{DM\mid L} = 0$ witnesses
saturation.
\end{definition}

A lawful allocation keeps the support, and a lawful acquisition is strict (\Cref{lem:d6}).


\subsection{The catalog}\label{sec:setting-catalog}

\Cref{tab:catalog} lists the sixteen entries with their status. The third column says what Lean
proves. The fourth names the instance of a law, or the certificate that is the open obligation of a
template. The field lists of all certificates are in \Cref{supp:fields}.

{\footnotesize
\begin{longtable}{@{}>{\raggedright\arraybackslash}p{3.3cm}>{\raggedright\arraybackslash}p{2.1cm}>{\raggedright\arraybackslash}p{4.3cm}>{\raggedright\arraybackslash}p{4.6cm}@{}}
\caption{The catalog.}\label{tab:catalog}\\
\toprule
\textbf{Entry} & \textbf{Status} & \textbf{Proved in Lean} & \textbf{Instance or open obligation} \\
\midrule
\endfirsthead
\toprule
\textbf{Entry} & \textbf{Status} & \textbf{Proved in Lean} & \textbf{Instance or open obligation} \\
\midrule
\endhead
\bottomrule
\endfoot
E1 internalization: a discharged recurring challenge has a carried repair source; perpetual
self-extension & law, weak instance & forced self-extension from recurrence and installation; clocked
device; five statuses exclusive & \Cref{prop:e1-forcing,thm:e1-clocked}; internalization status
remains a certificate (\Cref{tmpl:e1}) \\
E2 bounded reflexivity: audit towers are stratified and capacity-bounded & law, weak instance &
capacity bound over a declared inventory; two-level audited tower; priority status partition &
\Cref{thm:e2-declared,thm:e2-tower}; forced stratification remains a certificate \\
E3 self-maintaining reclosure: object closure and apparatus reinstatement are separate & law, weak
instance & executed repair with history-linked reinstatement & \Cref{thm:e3-episode};
maintained-closure status remains a certificate (\Cref{tmpl:e3}) \\
E4 repair compilation: repeated repair enters lower support & certificate specification &
compiled status from a full compilation case; brittleness from a split pair & compilation case
(\Cref{tmpl:e4}) \\
E5 reclosure collapse: no admissible carried rescue restores viability & definition & falsifier
refutes collapse; seven statuses exclusive & complete rescue inventory (\Cref{def:e5}) \\
E6 priced exposure allocation (attention) & general theorem & multiplier algebra; sufficiency for capped linear
objectives; optimum at a cap & numerical instance \Cref{thm:e6-atcap}; an instance on a carried
probe economy is open \\
E7 alarm: a blind spot that threatens viability is brought into current attention, declared an open
boundary, or discounted & conditional classification & the three dispositions exclusive & a unique
total disposition package for the specified blind spot and horizon (\Cref{tmpl:e7}) \\
E8 control price: carried control signals contain shadow-price components & definition & component
and summary lawfulness unfold definitions; seven statuses exclusive & component-lawful case
(\Cref{def:e8}) \\
E9 strict probe acquisition (curiosity) & general theorem & finite argmax; operator saturation gives
zero discharge & --- (\Cref{thm:e9}) \\
E10 cognitive demarcation, with intention and goal & conditional classification & statuses
exclusive for cognition, intention and goal & complete status records (\Cref{tmpl:e10}) \\
E11 institutional rewrite: a certified rewrite of what a lower layer can do & conditional
classification & statuses exclusive; a change of probabilities alone, or a label whose effect fades,
is not a rewrite & rewrite comparison (\Cref{tmpl:e11}) \\
E12 individuation: maximal jointly closed subcarrier with a sufficient boundary & law, weak instance & repair-closed
parts closed under intersection; individual containing its repair; overlapping maxima &
\Cref{thm:e12-repair,prop:e12-overlap}; budget and record closure remain a certificate
(\Cref{tmpl:e12}) \\
E13 repair transport: communication, teaching, coercion null, symbolic repair & law (paid transport), weak instance & composition; chain bound; transport across contexts & \Cref{thm:e13-chain,thm:e13-cross};
teaching, coercion and symbolic repair remain certificates (\Cref{tmpl:e13}) \\
E14 reconsolidation: a conflict exposed by recall is repaired, coarsened or recorded as unresolved & conditional
classification & lawful trichotomy and nine statuses exclusive & complete outcome inventory
(\Cref{tmpl:e14}) \\
E15 offline reclosure: debt discharge alternates in the deficit regime & conditional
classification & eight statuses exclusive; reallocation identity & complete status record
(\Cref{tmpl:e15}); identity in \Cref{prop:e15-identity} \\
E16 adaptability: routes that agree now but differ in future repair capacity & conditional
classification & eight statuses exclusive; agreement now and unequal later repair probabilities read
from the record & complete status record (\Cref{tmpl:e16}) \\
\end{longtable}
}
